\newtheorem{mtproposition}{Proposition}
\newtheorem{mttheorem}{Theorem}

\section{The original link function, including signed residuals}
Condition on a predictor fitted independently of the calibration and test data.
Let $\mathbf E^1,\ldots,\mathbf E^{n+1}\in\mathbb R^d$ be exchangeable
coordinatewise residual vectors, with the first $n$ observed. Write $N=n+1$,
$r=\lceil N(1-\alpha)\rceil$, and
\[
 Q_\alpha(v_1,\ldots,v_n)=v_{(r)}\quad(r\le n),
 \qquad Q_\alpha(v_1,\ldots,v_n)=+\infty\quad(r=N).
\]
Throughout, acceptance uses weak inequalities. Ties are permitted. Define
the calibration statistics $m_j=n^{-1}\sum_i E^i_j$ and
$s_j^2=n^{-1}\sum_i(E^i_j-m_j)^2$, initially with every $s_j>0$.
The original paper \cite{fan2025} augments these observations with a candidate residual
$z_j$ and uses
\begin{equation}\label{mt:augmented}
 \mu_j(z)=\frac{nm_j+z}{N},\qquad
 \sigma_j(z)=\sqrt{s_j^2+\frac{(z-m_j)^2}{N}},\qquad
 F_j(t,z)=\frac{t-\mu_j(z)}{\sigma_j(z)}.
\end{equation}
Here $s_j^2$ has divisor $n$, whereas the augmented sample variance has
divisor $N-1=n$. The distinction determines the constants below.
The candidate's self-score and scalarized calibration score are
\begin{equation}\label{mt:self}
 g_j(e)=F_j(e,e)=
 \frac{n(e-m_j)/N}{\sqrt{s_j^2+(e-m_j)^2/N}},\qquad
 \Phi_{\mathbf z}(\mathbf t)=\max_jF_j(t_j,z_j).
\end{equation}

The original nonnegative link has a positive middle branch
$m_j+Ns_j|c|/\sqrt{n^2-Nc^2}$, a negative middle branch
$\max\{0,m_j-Ns_j|c|/\sqrt{n^2-Nc^2}\}$, and outer values zero and
$+\infty$ at $c\le-n/\sqrt N$ and $c\ge n/\sqrt N$.
Its algebra already supports signed residuals: remove the truncation at zero,
preserve the sign of the middle root, and represent an empty lower branch by
$-\infty$. The resulting exact inverse is
\begin{equation}\label{mt:link}
 \omega_j(c)=\begin{cases}
 -\infty,&c\le -T,\\[2pt]
 m_j+\displaystyle\frac{Ns_jc}{\sqrt{n^2-Nc^2}},&-T<c<T,\\[5pt]
 +\infty,&c\ge T,
 \end{cases}\qquad T=\frac n{\sqrt N}.
\end{equation}
Indeed, $g_j'(e)=(n/N)s_j^2\{s_j^2+(e-m_j)^2/N\}^{-3/2}>0$,
its range is $(-T,T)$, and solving $g_j(e)=c$ requires
$\operatorname{sign}(e-m_j)=\operatorname{sign}(c)$. Hence, for finite $e$,
\begin{equation}\label{mt:inverse}
 g_j(e)\le c\quad\Longleftrightarrow\quad e\le\omega_j(c).
\end{equation}
For nonnegative scores one intersects this sublevel with $[0,\infty)$;
for signed scores one uses the appropriate attainable domain. Clamping a
negative endpoint to zero would turn an empty sublevel into a singleton, so
the unclamped form is useful even for a precise treatment of nonnegative ties.

At test input $x$, suppose $E_j(x,y_j)$ takes values in
$D_j(x)=[a_j(x),\infty)\cap\mathbb R$, allowing $a_j=-\infty$.
For absolute residuals, $E_j=|y_j-\hat f_j(x)|$ and $a_j=0$.
For signed CQR residuals \cite{romano2019}, with ordered fitted quantiles,
\begin{align}\label{mt:cqr}
 E_j(x,y_j)&=\max\{\hat q_j^-(x)-y_j,\ y_j-\hat q_j^+(x)\}
             =|y_j-\hat m_j(x)|-h_j(x),\notag\\
 \hat m_j(x)&=\tfrac12(\hat q_j^-+\hat q_j^+),\qquad
 h_j(x)=\tfrac12(\hat q_j^+-\hat q_j^-),\qquad a_j(x)=-h_j(x).
\end{align}
Thus a threshold $L_j$ gives $[\hat q_j^- - L_j,\hat q_j^+ + L_j]$;
negative thresholds shrink the fitted interval. The domain lower bound is
known from the fitted interval at $x$, and is not a minimum estimated from
the calibration scores. Those scores, observed at other inputs, may lie below
$a_j(x)$. The link is therefore not the obstacle to signed residuals; the
nonnegative scale separation in the old shortcut is.

\section{Review: GWC, full LWC, and the old shortcut}
\paragraph{Full conformal reference and GWC.}
For a candidate $\mathbf e=\mathbf E(x,y)$ define
\[
 q(\mathbf e)=Q_\alpha\bigl(\Phi_{\mathbf e}(\mathbf E^1),\ldots,
                                  \Phi_{\mathbf e}(\mathbf E^n)\bigr),\qquad
 C_{\rm full}(x)=\{y:\max_jg_j(e_j)\le q(\mathbf e)\}.
\]
This candidate-dependent set is the reference we enclose, and need not be
rectangular. For an interval $I$ put
$\psi_{j,I}(t)=\sup_{z\in I}F_j(t,z)$. GWC uses
\begin{equation}\label{mt:gwc}
 \Phi_G(\mathbf t)=\max_j\psi_{j,D_j(x)}(t_j),\quad
 q_G=Q_\alpha\bigl(\Phi_G(\mathbf E^i):i\le n\bigr),\quad
 U_j=\omega_j(q_G),\quad G_j=D_j(x)\cap(-\infty,U_j].
\end{equation}
Every candidate in $D(x)=\prod_jD_j(x)$ has $q(\mathbf e)\le q_G$;
hence $C_{\rm full}(x)\subseteq\{y:\mathbf E(x,y)\in G\}$, where
$G=\prod_jG_j$. This reviews the original GWC when $D_j=[0,\infty)$
and extends the same construction to signed scores.

\paragraph{The original full LWC.}
For nonnegative residuals, partition each $G_j=[0,U_j]$ at its interior
calibration residuals. We use closed cells
$I_{j,k}=[\eta_{j,k-1},\eta_{j,k}]\cap\mathbb R$ so shared boundary points
are retained. A full cell is $\mathcal I_{\mathbf h}=\prod_jI_{j,h_j}$;
there are at most $(n+1)^d$ such cells. The original paper uses
\begin{equation}\label{mt:separated}
 c_j=\inf_{z\in G_j}\frac{\mu_j(z)}{\sigma_j(z)},\qquad
 r_j(k)=\inf_{z\in I_{j,k}}\sigma_j(z),\qquad
 \Phi^{\rm sep}_{\mathbf h}(\mathbf t)
       =\max_j\left\{\frac{t_j}{r_j(h_j)}-c_j\right\},\quad\mathbf t\ge0.
\end{equation}
Its cell quantile $q^{\rm sep}_{\mathbf h}=Q_\alpha(
\Phi^{\rm sep}_{\mathbf h}(\mathbf E^i):i\le n)$ produces the local region
\[
 R^{\rm sep}_{\mathbf h}=
 \mathcal I_{\mathbf h}\cap\prod_j(-\infty,\omega_j(q^{\rm sep}_{\mathbf h})].
\]
Full LWC enumerates these regions, takes their union, and reports its
coordinatewise upper enclosure within $G$. This is a local bound on the
unknown candidate, followed by a rectangular enclosure; it is not full
conformal prediction itself.

For clarity, an \emph{exact coupled} full-cell alternative would replace
$\Phi^{\rm sep}_{\mathbf h}$ by
$\Phi^{\rm exact}_{\mathbf h}(\mathbf t)=
\max_j\psi_{j,I_{j,h_j}}(t_j)$.
It retains the whole ratio under each supremum and is no larger than the
separated cell score for nonnegative $\mathbf t$. It is a sharper conceptual
reference, not the separated full LWC defined in the original paper.
Neither full-cell construction is run in the meeting experiments.

\paragraph{The old row shortcut.}
The paper avoids enumerating all cells by varying one coordinate at a time,
holding the others in a cell containing the empirical mean. If GWC clips that
cell, the relevant point is $p_\ell=\operatorname{proj}_{G_\ell}(m_\ell)$,
with $\bar s_\ell=\sigma_\ell(p_\ell)$, the minimum scale on $G_\ell$.
The row score for coordinate $j$ and cell $k$ is
\begin{equation}\label{mt:oldrow}
 \Phi^{\rm old}_{j,k}(\mathbf t)=
 \max\left\{\frac{t_j}{r_j(k)}-c_j,
             \max_{\ell\ne j}\left(\frac{t_\ell}{\bar s_\ell}-c_\ell\right)
      \right\}.
\end{equation}
Its quantile, inverse link, and intersection with $I_{j,k}$ determine a
candidate upper endpoint; maximizing over $k$ gives $L_j^{\rm old}$.
The manuscript's row reduction applies under its mean-cell condition, and
the implemented shortcut falls back to GWC if this cell is unavailable.
The old shortcut here means this separated row construction with compatible
weak boundary conventions. Its key restriction is $t_\ell\ge0$:
$t_\ell/\sigma_\ell(z)\le t_\ell/\inf\sigma_\ell$ reverses for negative
$t_\ell$. Simply changing the inverse link cannot repair that inequality.

\section{The envelope construction}
\paragraph{An exact, inexpensive coordinate bound.}
Unlike scale separation, $\psi_{j,I}(t)$ is valid for either sign of $t$.
Suppressing $j$ and writing $I=[\ell,u]\cap\mathbb R$, differentiation gives
\begin{equation}\label{mt:derivative}
 \frac{\partial F(t,z)}{\partial z}
 =-\frac{s^2+(t-m)(z-m)}{N\{s^2+(z-m)^2/N\}^{3/2}}.
\end{equation}
For $t>m$, the sole stationary maximum is
$z^*(t)=m-s^2/(t-m)$, with value
$\sqrt{(t-m)^2/s^2+1/N}$. For $t<m$ the stationary point is a minimum,
and for $t=m$ the function is decreasing. Therefore
\begin{equation}\label{mt:supremum}
 \psi_{j,I}(t)=\max\left\{
 F_j(t,\ell),F_j(t,u),
 \left.\sqrt{(t-m_j)^2/s_j^2+1/N}\ \right|\ t>m_j,\ z^*(t)\in I
 \right\}.
\end{equation}
The last entry is omitted if inadmissible. Infinite endpoints use
$F_j(t,-\infty)=1/\sqrt N$ and $F_j(t,+\infty)=-1/\sqrt N$.
This formula computes all GWC and envelope scores without numerical optimization.

\paragraph{One coordinate cell, with the others bounded globally.}
Build GWC using \eqref{mt:gwc}, now with either residual domain. Subdivide
each nonempty $G_j$ at its distinct interior calibration values, retaining
singletons when necessary. For each $j$, precompute
\begin{align}\label{mt:surface}
 H_j(\mathbf t)&=\max_{\ell\ne j}\psi_{\ell,G_\ell}(t_\ell),\notag\\
 \overline\Phi_{j,k}(\mathbf t)
  &=\max\{\psi_{j,I_{j,k}}(t_j),H_j(\mathbf t)\}
   =\sup_{\substack{z_j\in I_{j,k}\\z_\ell\in G_\ell,\ \ell\ne j}}
      \Phi_{\mathbf z}(\mathbf t),\notag\\
 \overline q_{j,k}&=Q_\alpha\bigl(\overline\Phi_{j,k}(\mathbf E^i):i\le n\bigr).
\end{align}
For $d=1$, set $H_1=-\infty$. Thus the envelope keeps one coordinate local
and lets all other candidate coordinates range over GWC. It retains the
coupled location--scale ratio while replacing exponential full-cell
enumeration by coordinate searches. Define
\begin{align}\label{mt:pieces}
 A_{j,k}&=I_{j,k}\cap(-\infty,\omega_j(\overline q_{j,k})],\qquad
 L_j=\max_k B_{j,k},\notag\\
 B_{j,k}&=\begin{cases}
 \min\{\eta_{j,k},\omega_j(\overline q_{j,k})\},&A_{j,k}\ne\varnothing,\\
 -\infty,&A_{j,k}=\varnothing,
 \end{cases}
\end{align}
The final region is
$C_{\rm env}(x)=\{y:E_j(x,y_j)\in D_j(x),\ E_j(x,y_j)\le L_j\ \forall j\}$.
It is empty if any coordinate has no accepted piece. It can fill gaps between
pieces, and generally differs from both full conformal and full-cell LWC.
For ordinary absolute or CQR residuals, its outcome sublevels are intervals,
so the final region is a box.

\paragraph{Search and cost.}
With the calibration coordinates sorted, scan cells from right to left and
stop at the first nonempty piece. This is exact: every earlier cell ends at
or below the surviving cell's left endpoint, and therefore cannot give a
larger $B_{j,k}$. No unproved monotonicity of acceptance is required.
All off-coordinate values $H_j(\mathbf E^i)$ can be precomputed in $O(nd)$
using prefix and suffix maxima. Each inspected cell then costs $O(n)$ with
order-statistic selection. Total worst-case cost is
$O(dn\log n+dn^2)$, and actual search cost is $O(n\sum_j M_j)$ when
$M_j$ cells are inspected. This establishes tractability, not uniform runtime
dominance. Calibration statistics can be reused; input-dependent CQR widths
may require different domain calculations at different test inputs.

\section{Validity and uniform improvement: separate proofs}
\begin{mtproposition}[Finite-sample validity]\label{mt:validity}
For every realized positive-scale calibration sample and test input,
$C_{\rm full}(x)\subseteq C_{\rm env}(x)\subseteq C_{\rm GWC}(x)$.
Under exchangeability, $\Pr\{Y^{n+1}\in C_{\rm env}(X^{n+1})\}\ge1-\alpha$.
\end{mtproposition}
\begin{proof}
At the true test residual, the augmented means and scales are symmetric
functions of all $N$ observations. Applying the same standardized maximum
to all observations gives exchangeable scalar scores. The conformal rank
argument, including weak acceptance for ties, gives coverage at least
$r/N\ge1-\alpha$ for $C_{\rm full}$.

For an accepted candidate $\mathbf e$, GWC containment gives $\mathbf e\in G$.
Choose for each $j$ a covering cell $I_{j,k}$ containing $e_j$.
Every calibration score obeys
$\Phi_{\mathbf e}(\mathbf E^i)\le\overline\Phi_{j,k}(\mathbf E^i)$.
Coordinatewise ordering of a score vector preserves ordering of its $r$th
order statistic; hence $q(\mathbf e)\le\overline q_{j,k}$.
Full acceptance and the increasing link imply
$e_j\le\omega_j(\overline q_{j,k})$, so $e_j\in A_{j,k}$ and $e_j\le L_j$.
This holds for every coordinate without a Bonferroni adjustment.
Finally, every piece lies in $G_j$, giving $L_j\le U_j$.
\end{proof}

\begin{mttheorem}[Uniform weak improvement over the old shortcut]\label{mt:dominance}
Use the same nonnegative residuals, calibration sample, $\alpha$, GWC domain,
and closed cells for both methods, with $s_j>0$. Then
\[
 C_{\rm env}(x)\subseteq C_{\rm old}(x).
\]
For nonempty regions this means $L_j\le L_j^{\rm old}$ in every coordinate,
and therefore weakly smaller outcome intervals and volume. Envelope coverage
remains at least $1-\alpha$.
\end{mttheorem}
\begin{proof}
For $t_j\ge0$ and $z\in I_{j,k}$,
\[
 F_j(t_j,z)=\frac{t_j}{\sigma_j(z)}-\frac{\mu_j(z)}{\sigma_j(z)}
 \le\frac{t_j}{r_j(k)}-c_j.
\]
For an off-coordinate $\ell$, the same inequality on $G_\ell$ gives
$\psi_{\ell,G_\ell}(t_\ell)\le t_\ell/\bar s_\ell-c_\ell$.
Consequently $\overline\Phi_{j,k}(\mathbf E^i)\le
\Phi^{\rm old}_{j,k}(\mathbf E^i)$ for every $i,j,k$.
Quantile monotonicity and \eqref{mt:inverse} imply that every envelope piece
is contained in the corresponding old piece. Taking upper endpoints and
then maxima proves the result. If the old shortcut falls back to GWC,
Proposition~\ref{mt:validity} gives containment directly.
\end{proof}

The separation bound can be strict because the smallest scale and smallest
location--scale ratio occur at different candidate values. For a concrete
one-dimensional example, take calibration residuals $(1,2,3)$ and
$\alpha=1/2$. Then $m=2$, $s^2=2/3$, $r=2$, and GWC ends at
$U=2+2/\sqrt{21}$. On $[2,U]$, the envelope's median score is zero,
so $L^{\rm env}=2$; the separated median is
$\sqrt6-\sqrt{27/20}$, whose link exceeds $U$, so $L^{\rm old}=U$.
The absolute-residual interval is strictly shorter. Equality is also
possible, for example when $r=N$ and both methods return the whole domain.
The theorem does not assert strict improvement on every sample, higher
coverage than a superset, or dominance over exact coupled full-cell LWC.

\paragraph{Scope of the signed comparison.}
Capping $E$ at zero loses negative values; shifting by a width that changes
with $x$ changes relative residuals across observations. The preceding theorem
does not compare those representations. A fixed coordinatewise shift has
additional structure: if a training-only constant $b_j$ satisfies
$E_j(x,y)\ge-b_j$ everywhere, then
$F_j^{\rm shift}(t+b_j,z+b_j)=F_j(t,z)$ and
$\omega_j^{\rm shift}(c)=\omega_j(c)+b_j$.
The signed attainable domain translates into a subset of $[0,\infty)$.
Taking suprema on that smaller domain and applying the same proof shows that
the signed envelope is contained in the old constant-shifted region, under
matching exact boundary rules. No analogous containment is claimed against
capping or input-dependent width shifts; their practical comparison is empirical.

\paragraph{Boundary and implementation conventions.}
If $r=N$, return the whole attainable domain. If any observed $s_j=0$, the
experiments use a common whole-domain fallback; the displayed dominance proof
concerns positive scales. A valid full-conformal reference on degenerate
samples uses augmented scale $1$ when that entire augmented coordinate is
constant; this symmetric rule and the whole-domain fallback preserve coverage.
Closed cells retain singleton endpoints. Mathematical inequalities refer to
exact arithmetic with these conventions; the old executable's numerical
jitter and finite-precision outputs are checked separately. All coverage
statements average over calibration and test data, rather than asserting
coverage conditional on a particular input or calibration sample.
